\documentclass[xcolor=table,aspectratio=169]{beamer}
\usepackage{beamerthemesplit}
\usepackage{wrapfig}
\usetheme{SPbGU}
\usepackage{pdfpages}
\usepackage{amsmath}
\usepackage{cmap}
\usepackage[T2A]{fontenc}
\usepackage[utf8]{inputenc}
\usepackage[english]{babel}
\usepackage{indentfirst}
\usepackage{amsmath}
\usepackage{tikz}
\usepackage{multirow}
\usepackage[noend]{algpseudocode}
\usepackage{algorithm}
\usepackage{algorithmicx}
\usepackage{fancyvrb}
\usepackage{fontawesome5}
\usepackage{fontawesome}

\usetikzlibrary{arrows, shapes, shapes.geometric, fit, automata, positioning, calc}
\usetikzlibrary{decorations.pathmorphing}
\tikzset{state/.style={draw,circle}}
\tikzset{state none/.style={draw,none}}

\newcommand{\tikzmark}[1]{\tikz[overlay,remember picture] \node (#1) {};}
\def\Put(#1,#2)#3{\leavevmode\makebox(0,0){\put(#1,#2){#3}}}

\beamertemplatenavigationsymbolsempty

\title[CFPQ: Obstacles on the Way to Adoption]{Context-Free Path Querying: Obstacles on the Way to Adoption}
\institute[JBR, SPbSU]{
JetBrains Research, Programming Languages and Tools Lab  \\
St. Petersburg State University\\
\vspace{0.7cm}
\url{https://research.jetbrains.org/groups/plt_lab/}
}

\author[Semyon Grigorev]{Semyon Grigorev}

\date{16.07.2021}

\begin{document}
{
\begin{frame}[fragile]
  \begin{tabular}{p{2.0cm} p{9.5cm} p{1cm}}
   \begin{center}
      \includegraphics[height=1.5cm]{pictures/jetbrainsResearch.pdf}
    \end{center}
    &
    \begin{center}
      LDBC TUC 
    \end{center}
    &
    \begin{center}
      \includegraphics[height=1.5cm]{pictures/SPbGU_Logo.png}
    \end{center}
  \end{tabular}
  \titlepage
\end{frame}
}

%\begin{frame}[fragile]
%
%  \frametitle{Who we are?}
%
%  \begin{minipage}[t]{0.48\textwidth}
%    \begin{itemize}
%    \item !!!
%    \end{itemize}    
%  \end{minipage}
%  \begin{minipage}[t]{0.48\textwidth}
%    \begin{itemize}
%    \item !!!
%    \end{itemize}    
%  \end{minipage}
%
%\end{frame}

\begin{frame}[fragile]
 
  \frametitle{Formal Language Constrained Path Querying}
  Navigational queries in edge-labelled graph

  \begin{minipage}[t]{0.43\textwidth}    
    \begin{center}
      \only<1>{
    \begin{tikzpicture}[shorten >=1pt,auto]
      \node[state] (q_0)                        {$0$};
      \node[state] (q_1) [right=of q_0]         {$1$};
      \node[state] (q_2) [below=of q_1]    {$2$};
      \node[state] (q_3) [below=of q_0]    {$3$};
      \node[state] (q_5) [right=of q_1]   {$5$};
      \node[state] (q_4) [below=of q_5]   {$4$};
      \path[->]
       (q_0) edge  node {$a$} (q_1)
       (q_1) edge  node {$b$} (q_5)
       (q_1) edge[left]  node {$b$} (q_2)
       (q_3) edge[above]  node {$c$} (q_2)
       (q_4) edge[bend left]  node {$c$} (q_3)
       (q_2) edge[above]  node {$c$} (q_4)
       (q_5) edge[bend left, right]  node {$d$} (q_4)
       (q_4) edge[bend left, left]  node {$d$} (q_5);
   \end{tikzpicture}}
   \only<2>{
   \begin{tikzpicture}[shorten >=1pt,auto]
    \node[state] (q_0)                        {$0$};
    \node[state] (q_1) [right=of q_0]         {$1$};
    \node[state] (q_2) [below=of q_1]    {$2$};
    \node[state] (q_3) [below=of q_0]    {$3$};
    \node[state] (q_5) [right=of q_1]   {$5$};
    \node[state] (q_4) [below=of q_5]   {$4$};
    \path[->]
     (q_0) edge [red] node {$a$} (q_1)
     (q_1) edge [red] node {$b$} (q_5)
     (q_1) edge[left]  node {$b$} (q_2)
     (q_3) edge[above]  node {$c$} (q_2)
     (q_4) edge[bend left, red]  node {$c$} (q_3)
     (q_2) edge[above]  node {$c$} (q_4)
     (q_5) edge[bend left, right, red]  node {$d$} (q_4)
     (q_4) edge[bend left, left]  node {$d$} (q_5);
 \end{tikzpicture}}
 \only<3->{
   \begin{tikzpicture}[shorten >=1pt,auto]
    \node[state] (q_0)                        {$0$};
    \node[state] (q_1) [right=of q_0]         {$1$};
    \node[state] (q_2) [below=of q_1]    {$2$};
    \node[state] (q_3) [below=of q_0]    {$3$};
    \node[state] (q_5) [right=of q_1]   {$5$};
    \node[state] (q_4) [below=of q_5]   {$4$};
    \path[->]
     (q_0) edge [blue] node {$a$} (q_1)
     (q_1) edge node {$b$} (q_5)
     (q_1) edge[left, blue]  node {$b$} (q_2)
     (q_3) edge[above]  node {$c$} (q_2)
     (q_4) edge[bend left]  node {$c$} (q_3)
     (q_2) edge[above, blue]  node {$c$} (q_4)
     (q_5) edge[bend left, right]  node {$d$} (q_4)
     (q_4) edge[bend left, left, blue]  node {$d$} (q_5);
 \end{tikzpicture}}
  \end{center}
  \begin{itemize}
    \item Path to find: $$0 \xrightarrow{a} v_0 \xrightarrow{b} v_1 \underbrace{\xrightarrow{d} v_2 \xrightarrow{c} v_3 \ldots v_k \xrightarrow{c}}_{c \text{ or } d \text{ in arbitrary order}} v$$
    \pause     %\item Respective regular constraint: $a b (c \mid d)^*$
    \end{itemize} 
  \end{minipage}  
  \pause \pause
  \begin{minipage}[t]{0.56\textwidth}
    \vspace{-3.5cm}
    \begin{itemize}
      \item $w(v_0 \xrightarrow{l_0} v_1 \xrightarrow{l_1} \ldots \xrightarrow{l_{k-1}} v_k) = l_0 l_1 \ldots l_{k-1}$
      \item $Q = \{(v_i,v_j) \mid \exists \pi = v_i \to \ldots \to v_j; w(\pi) \in \mathcal{L}\}$, where $\mathcal{L}$ --- formal language         
    \begin{itemize}
    \item[\faCheck] Regular, RPQ ($ab(c \mid d)^*$)
    \item[\faGears] \textbf{Context-Free}, CFPQ ($a^n b^n$)
    \item[\faHourglassHalf] Multiple Context-Free ($a^nc^mb^nd^m$)
    \end{itemize}    
  \end{itemize} 
  \pause
  Variations:
  \begin{itemize}
    \item All-pairs
    \item Multiple source
    \item Reachability
    \item All paths
    \item \ldots
    \end{itemize}    
  \end{minipage}

\end{frame}


\begin{frame}[fragile]
  \frametitle{Applications of Context-Free Path Querying}
  \begin{center}
    \fbox{
  \parbox{0.95\linewidth}
  {
    \begin{center}
  Hierarchy analysis: variations of the \textit{same-generation queries} is the essence of CFPQ
\end{center}
}
}
\pause
\fbox{
  \begin{minipage}[t][3.5cm][t]{0.46\textwidth}
    Graph structured data analysis
    \begin{itemize}
      \item Introduced by M. Yannakakis in 1990\footnotemark
      \item Biological data analysis
      \item Data provenance analysis
      \item \ldots
    \end{itemize}
  \end{minipage}
  }
  \pause
  \fbox{
  \begin{minipage}[t][3.5cm][t]{0.46\textwidth} 
    Static code analysis
    \begin{itemize}
      \item Introduced by T. Reps in 1997\footnotemark
      \item Interprocedural points-to analysis
      \item Interprocedural alias analysis
      \item Type inference related tasks
      \item \ldots 
    \end{itemize}
  \end{minipage}}
  
  \addtocounter{footnote}{-1}
  \footnotetext{Thomas Reps. 1997. ``Program Analysis via Graph Reachability''.}
  \addtocounter{footnote}{1}
  \footnotetext{Mihalis Yannakakis. 1990. ``Graph-theoretic Methods in Database Theory''.}
  \pause
  \fbox{
  \parbox{0.95\linewidth}
  {
    \begin{center}
    Graph databases
  \end{center}
  }
  }
\end{center}
\end{frame}

\begin{frame}[fragile]
  \frametitle{There is No Unified Infrastructure to Compare CFPQ Solutions}  
  \begin{itemize}
    \item[\faQuestion] Which algorithm is better for the specific problem?    
    \item[\faQuestion] How to assess if a newly developed algorithm is better than the existing ones?
    \item[\faQuestion] \ldots 
    \pause 
    \item[\faLightbulb] \textbf{Unified benchmarks}: dataset, workflow, \ldots
    \pause 
    \item[\faExclamation] The only attempt to compare different algorithms: ``An Experimental Study of Context-Free Path Query Evaluation Methods''\footnote{Jochem Kuijpers, George Fletcher, Nikolay Yakovets, and Tobias Lindaaker. 2019}
      \pause
      \begin{itemize}
      \item[\faFrown] ``We conclude that state of the art solutions are not able to cope with large graphs as found in practice.''
      \end{itemize}
    \end{itemize}
  \pause
  \vfill
  Difficulties
    \begin{itemize}
      \item Data is spread over projects and papers in different communities
      \item There is a huge number of different subclasses of the problem
      \begin{itemize}
        \item all-pairs, single source, multiple source, \dots 
        \item reachability, single path, all path, \dots 
      \end{itemize}
  \end{itemize}
\end{frame}

\begin{frame}[fragile]
  \frametitle{There is No Support of CFPQ in Real-World Graph Databases}
  \begin{itemize}        
    \item[\faQuestion] Which database or graph analysis system should you choose?  
    \begin{itemize}
      \item[\faFrown] H. Miao and A. Deshpande: ``Though the problem has been first studied in our community [40], there is little follow up and support in the context of modern graph databases \dots''\footnote{H. Miao and A. Deshpande,  ``Understanding Data Science Lifecycle Provenance via Graph Segmentation and Summarization'', 2019}       
    \end{itemize}
  \end{itemize}
  \vfill
  Difficulties
  \begin{itemize}        
    \item[\faQuestion] How to choose an appropriate algorithm for query engine?
    \begin{itemize}
      \item[\faGears] Benchmarks for querying algorithms
    \end{itemize}
    \item[\faQuestion] How to express context-free constraints in graph query language? 
    \begin{itemize}
      \item[\faQuestion] Syntactic features to express context-free language constraints
      \item[\faQuestion] Semantics of quey language
      \pause 
      \item[\faGears] GQL standard  
    \end{itemize}
  \end{itemize}  
\end{frame}


\begin{frame}[fragile]
  \frametitle{Our Results}

  \begin{itemize}
    \item[\faCheck] Collection of linear algebra based algorithms for CFPQ
    \begin{itemize}
      \item SuiteSparse is utilized for sparse linear algebra subroutines
      \item Published: \url{https://github.com/JetBrains-Research/CFPQ_PyAlgo}
    \end{itemize}
    \pause
    % Definitely, not all abilities of Cypher language is supported, but ...
    \item[\faCheck] Full-stack support of CFPQ\footnote{Arseniy Terekhov et al. 2019. ``Multiple-Source Context-Free Path Querying in Terms of Linear Algebra''.}    
    \begin{itemize}
      \item On top of RedisGraph: query engine is extended with CFPQ algorithm
      \item openCypher is extended to support CFPQ
    \end{itemize} 
    \pause 
    \item[\faGears] Dataset for CFPQ benchmarking: early stages 
    \begin{itemize}
      \item Synthetic graphs
      \begin{itemize}
        \item Theoretical worst case
        \item Complicated cases
      \end{itemize}
      \item Real-world graphs
      \begin{itemize}
        \item Static code analysis
        \item Biological data analysis
      \end{itemize}
      \item Published: \url{https://github.com/JetBrains-Research/CFPQ_Data} 
    \end{itemize}     
  \end{itemize}

\end{frame}

\begin{frame}[fragile]
  \frametitle{Our Results Evaluation}
  \begin{minipage}[t]{0.48\textwidth}
    \begin{itemize}
      \item All-pairs reachability queries
      \item \textit{geospecies, taxonomy} --- biological data
      \item \textit{crypto, drivers, fs} --- points-to analysis      
      \item Time in seconds      
    \end{itemize}
  \end{minipage}
  \begin{minipage}[t]{0.48\textwidth}
    \begin{itemize}
      \item GPU: Geforce GTX 1070, 1.5GHz, 8Gb RAM, 1920 CUDA cores 
      \item CPU: Intel core i7-6700 CPU, 3.4GHz, DDR4 64Gb RAM
    \end{itemize}
  \end{minipage}
\begin{center}
  \rowcolors{2}{black!2}{black!10}
  \tikzmark{zzz}{ }
  \begin{tabular}{| c | c | c || c | c | c | c |}
    \hline
    Graph      & \#V       & \#E         & Neo4j\footnote{Jochem Kuijpers et al. 2019. ``An Experimental Study of Context-Free Path Query Evaluation Methods''.} 
                                                               & RedisGraph\footnote{Arseniy Terekhov et al. 2019. ``Multiple-Source Context-Free Path Querying in Terms of Linear Algebra''.} 
                                                                                          & Lin.al. CPU\footnote{Standalone linear algebra based algorithm on CPU, using SuiteSparse.} 
                                                                                                                      & Lin.al. GPU\footnote{Standalone linear algebra based algorithm on GPU, using spbla.} \\
    \hline
    \hline
    geospecies & 450 609   & 2 311 461   & 6 953.9             & 80.1                     & 7.1                       & 0.8                         \\ 
    taxonomy   & 5 728 398 & 14 922 125  & n.a.                & \faGears                 & 1.1                       & 0.7                         \\
    %arch       & 3 448 422 & 5 940 484   & n.a.                & \faGears                 & 83.7                      &  27.9                       \\
    crypto     & 3 464 970 & 5 976 774   & n.a.                & \faGears                 & 84.8                      &  28.1                       \\
    drivers    & 4 273 803 & 7 415 538   & n.a.                & \faGears                 & 269.9                     & 62.5                        \\
    fs         & 4 177 416 & 7 218 746   & n.a.                & \faGears                 & 165.1                     & 47.7                        \\
    \hline
  \end{tabular}  
\end{center}
\pause
\onslide<2>{\tikz[overlay,remember picture]{\draw[draw=red,thick,double,fill opacity=0.2] ($ (zzz) + (-7.6,-1.1)$) rectangle ($ (zzz) + (7.9,-0.6)$);}}
\onslide<3>{\tikz[overlay,remember picture]{\draw[draw=red,thick,double,fill opacity=0.2] ($ (zzz) + (0.45,-0.6)$) rectangle ($ (zzz) + (2.85,-3.0)$);}}
\end{frame}


\begin{frame}[fragile]

  \frametitle{Ongoing Research}
\begin{itemize}
      \item[\faGears] Benchmarking of linear algebra based algorithms
      \begin{itemize}
        \item Comparison of different algorithms for different query semantics
        \item Investigation of scalability on multicore machines
        \item Estimation of performance on GPGPU
      \end{itemize}
      \item[\faGears] Developing and evaluating GLL-based CFPQ algorithm for Neo4j
      \begin{itemize}
        \item Multiple-source
        \item All paths and reachability-only
      \end{itemize}
      \pause
      \item[\faHourglassHalf] Describing semantics of (subset of) openCypher in terms of linear algebra (in Coq)
      \item[\faHourglassHalf] Utilizing multiple context-free languages as path constraints      
\end{itemize}
\end{frame}

\begin{frame}[fragile]
  \frametitle{What Should We Do?}
  \begin{itemize}
    \item[\faRoad] Publish unified benchmarks for formal language constrained path querying algorithms
    \begin{itemize}
      \item Graphs: synthetic and real-world
      \item Queries: templates and real-world queries
      \item Tasks: all-pairs, single source, reachability, \ldots            
    \end{itemize}
    \pause
    \item[\faRoad] Organize language constrained path querying competition
    \begin{itemize}
      \item Part of existing competition for graph processing systems
      \item Involve static analysis community
    \end{itemize}
    \pause
    \item[\faQuestion] Provide graph database support
    \begin{itemize}
      \item Different algorithms for different systems
      \item Syntax and semantics of query languages 
    \end{itemize}
  \end{itemize}
\end{frame}

\end{document}

